\begin{answer}
    
    We want to show that $\ell_{\text{semi-sup}}(\theta^{(t+1)}) \geq \ell_{\text{semi-sup}}(\theta^{(t)})$:

    For this we can use the fact from the standard derivation of EM, that for any $\Theta$, the log-likelihood of the unsupervised data is greater than or equal to the lower bound, we called in the E-Step:

    1) Let's start by looking at the difference between the objective function at step $t+1$ and step $t$.

    $$\ell_{\text{semi-sup}}(\theta^{(t+1)}) - \ell_{\text{semi-sup}}(\theta^{(t)})$$

    $$= \left( \ell_{\text{unsup}}(\theta^{(t+1)}) + \alpha \ell_{\text{sup}}(\theta^{(t+1)}) \right) - \left( \ell_{\text{unsup}}(\theta^{(t)}) + \alpha \ell_{\text{sup}}(\theta^{(t)}) \right)$$

    $$= \left( \ell_{\text{unsup}}(\theta^{(t+1)}) - \ell_{\text{unsup}}(\theta^{(t)}) \right) + \alpha \left( \ell_{\text{sup}}(\theta^{(t+1)}) - \ell_{\text{sup}}(\theta^{(t)}) \right)$$

    2) From the derivation of the standard EM algorithm, we know that the change in the unsupervised log-likelihood is lower-bounded by the change in its auxiliary function. For any choice of distributions. For any choice of $Q_i$, we have:

    $$\ell_{\text{unsup}}(\theta) \ge \sum_{i=1}^{n} \sum_{z^{(i)}} Q_i(z^{(i)}) \log \frac{p(x^{(i)}, z^{(i)}; \theta)}{Q_i(z^{(i)})}$$

    In the E-step, we set $Q_i^{(t)}(z^{(i)}) = p(z^{(i)} | x^{(i)}; \theta^{(t)})$. With this choice the lower bound is at $\theta = \theta^{(t)}$. This gives us with the key inequality from the Standard EM proof:

    \begin{align}
    \ell_{\text{unsup}}(\theta^{(t+1)}) - \ell_{\text{unsup}}(\theta^{(t)}) &\geq \left( \sum_{i=1}^{n} \sum_{z^{(i)}} Q_i^{(t)}(z^{(i)}) \log \frac{p(x^{(i)}, z^{(i)}; \theta^{(t+1)})}{Q_i^{(t)}(z^{(i)})} \right) \\
    &\quad - \left( \sum_{i=1}^{n} \sum_{z^{(i)}} Q_i^{(t)}(z^{(i)}) \log \frac{p(x^{(i)}, z^{(i)}; \theta^{(t)})}{Q_i^{(t)}(z^{(i)})} \right)
    \end{align}

    3) Now we can substitute this inequality into the expression from Step 1. For this we can define $J(\theta, Q^{(t)}) = \sum_{i=1}^{n} \sum_{z^{(i)}} Q_i^{(t)}(z^{(i)}) \log \frac{p(x^{(i)}, z^{(i)}; \theta)}{Q_i^{(t)}(z^{(i)})}$: 

    The inequality becomes:

    $$\ell_{\text{semi-sup}}(\theta^{(t+1)}) - \ell_{\text{semi-sup}}(\theta^{(t)}) \geq \left( J(\theta^{(t+1)}, Q^{(t)}) - J(\theta^{(t)}, Q^{(t)}) \right) + \alpha \left( \ell_{\text{sup}}(\theta^{(t+1)}) - \ell_{\text{sup}}(\theta^{(t)}) \right)$$

    The right-hand side becomes:

    $$\geq \left[ J(\theta^{(t+1)}, Q^{(t)}) + \alpha \ell_{\text{sup}}(\theta^{(t+1)}) \right] - \left[ J(\theta^{(t)}, Q^{(t)}) + \alpha \ell_{\text{sup}}(\theta^{(t)}) \right]$$

    4) The terms in the square brackets are now exactly the functions that we try to maximize in the M-step. As problem statement, it is defined as:

    $$\theta^{(t+1)} := \arg\max_{\theta} \left[ J(\theta, Q^{(t)}) + \alpha \ell_{\text{sup}}(\theta) \right]$$

    And this implies that the function at $\theta^{(t+1)}$ needs to be greater than or equal to the value at any other $\theta$ (including $\theta^{(t)}$). Hence, we can say:

    $$J(\theta^{(t+1)}, Q^{(t)}) + \alpha \ell_{\text{sup}}(\theta^{(t+1)}) \geq J(\theta^{(t)}, Q^{(t)}) + \alpha \ell_{\text{sup}}(\theta^{(t)})$$

    With this, we have shown that the objective function increases, as the right-hand side is always non-negative. This shows that $\ell_{\text{semi-sup}}(\theta^{(t+1)}) - \ell_{\text{semi-sup}}(\theta^{(t)}) \geq 0$
\end{answer}
